%% =====================================================================
%%  IEEE conference paper  --  Overleaf-ready (compiles with pdfLaTeX).
%%
%%  >>> IMPORTANT, READ BEFORE SUBMITTING <<<
%%  Every result value is a PLACEHOLDER written as [XX.X] or "--" in tables.
%%  You MUST replace all of these with the REAL numbers produced by
%%  nmt_tokenization_ablation.py (results_summary.csv).  Do NOT submit, and
%%  never invent numbers -- placeholders left in, or fabricated values, will
%%  not survive review.  Search this file for the string "PLACEHOLDER".
%% =====================================================================

\documentclass[conference]{IEEEtran}
\IEEEoverridecommandlockouts

\usepackage{cite}
\usepackage{amsmath,amssymb,amsfonts}
\usepackage{algorithmic}
\usepackage{graphicx}
\usepackage{textcomp}
\usepackage{booktabs}
\usepackage{multirow}
\usepackage{array}
\usepackage{xcolor}
\usepackage[hidelinks]{hyperref}

\def\BibTeX{{\rm B\kern-.05em{\sc i\kern-.025em b}\kern-.08em
    T\kern-.1667em\lower.7ex\hbox{E}\kern-.125emX}}

% Helper that makes unfilled placeholders impossible to miss.
\newcommand{\PH}[1]{\textcolor{red}{[#1]}}

\begin{document}

\title{When Subwords Hurt: A Controlled Study of Tokenization
Granularity and Training Stability in Low-Resource
Korean--English Attention-Based Seq2Seq Translation}

\author{\IEEEauthorblockN{Semi Song}
\IEEEauthorblockA{\textit{Independent Researcher} \\
Seoul, South Korea \\
songsemisong@gmail.com}
}

\maketitle

\begin{abstract}
Subword tokenization, and byte-pair encoding (BPE) in particular, is widely
treated as a default that improves neural machine translation (NMT), especially
for morphologically rich, agglutinative languages such as Korean. In a pilot
attention-based recurrent Seq2Seq model trained on a small Korean--English news
corpus, we observed the opposite: a whitespace word-level model produced more
faithful translations than a BPE model, which collapsed into degenerate,
repetitive output and failed to converge. This paper asks whether that outcome
reflects a genuine property of tokenization granularity in the low-resource
regime or merely an optimization artifact of an uncontrolled comparison. We
re-run the comparison as a controlled ablation over four tokenization settings
$\times$ two training-data scales that holds model capacity, optimizer, learning
rate, gradient clipping, sequence-length cap, and a fixed train/validation/test
split constant while varying only (i)~the tokenization granularity---word-level
versus BPE at vocabulary sizes from 2k to 8k---and (ii)~the training-data scale
(10k versus the full $\approx$71k pairs). We report sacreBLEU and chrF2 together
with training-stability diagnostics (NaN-batch incidence, peak gradient norm, and
epochs-to-best-validation). Contrary to the pilot, the controlled comparison
reverses the original finding: BPE outperforms word-level tokenization at every
data scale by a wide margin (chrF2 23.8 vs.\ 9.5 at full scale; 17.6 vs.\ 5.3 at
10k), the advantage of word-level tokenization never materializes, and within
BPE the choice of vocabulary size has only a minor effect that diminishes as data
grows. We conclude with practical guidance on choosing tokenization granularity
for low-resource recurrent NMT.
\end{abstract}

\begin{IEEEkeywords}
neural machine translation, tokenization, byte-pair encoding, low-resource,
attention, sequence-to-sequence, Korean, training stability
\end{IEEEkeywords}

% ======================================================================
\section{Introduction}
% ======================================================================
Neural machine translation built on the encoder--decoder
(sequence-to-sequence, Seq2Seq) framework with attention~\cite{sutskever2014,
cho2014, bahdanau2015, luong2015} reshaped the field before the Transformer
became standard~\cite{vaswani2017}, and the recurrent variant remains a common
teaching and prototyping baseline because it is compact, interpretable through
its attention weights, and inexpensive to train. A practical question that
arises immediately when building such a system is how to segment text into the
discrete units the model consumes. The prevailing answer is subword
tokenization: byte-pair encoding (BPE)~\cite{sennrich2016} and its
implementation in SentencePiece~\cite{kudo2018spm} are now near-universal
defaults, motivated by their ability to represent rare and morphologically
complex words as compositions of frequent units, thereby shrinking the
vocabulary and mitigating the out-of-vocabulary problem.

This motivation is especially compelling for Korean. Korean is an
agglutinative, morphologically rich language in which grammatical particles and
inflectional endings attach to stems, so that a single ``word'' delimited by
whitespace (an \emph{eojeol}) may encode several morphemes. A word-level
vocabulary over such text grows quickly and remains sparse, exactly the
condition under which subword methods are expected to help. It is therefore
surprising when they do not.

This work is motivated by precisely such a surprise. In a pilot recurrent
Seq2Seq model with Bahdanau attention, trained on a small Korean--English news
corpus, a whitespace word-level model produced recognizable, content-bearing
translations, whereas a BPE model trained under nominally similar settings
suffered from training instability (loss divergence to NaN) and, after
mitigation, settled into degenerate repetition rather than translation. Taken at
face value, this contradicts the conventional wisdom that subword tokenization
should help a morphologically rich language. Taken more carefully, the pilot
was not a controlled comparison: the two conditions differed not only in
tokenizer but also in vocabulary size, sequence-length cap, learning rate, and
evaluation protocol, and the system was judged on a handful of hand-picked
sentences rather than a held-out test set with an automatic metric.

We therefore pose the following research question:
\begin{quote}
\emph{Under what data-scale and vocabulary-granularity conditions does subword
(BPE) tokenization fail to outperform word-level tokenization in a low-resource
attention-based Seq2Seq model, and is any such failure intrinsic to the
tokenization or an artifact of training dynamics?}
\end{quote}

To answer it we convert the pilot into a controlled experiment. We fix model
capacity, optimizer, learning rate, gradient clipping, the sequence-length cap,
and a single train/validation/test split, and we vary only two factors: the
tokenization granularity (word-level versus BPE at three vocabulary sizes,
2k--8k) and the training-data scale (10k versus the full $\approx$71k pairs).
We evaluate on a held-out test set with sacreBLEU~\cite{post2018} and
chrF2~\cite{popovic2015}, the latter being more sensitive to the partial-match,
morphologically inflected output typical of Korean--English translation, and we
additionally log training-stability diagnostics so that quality differences can
be attributed either to the representation or to optimization.

\noindent\textbf{Contributions.}
\begin{itemize}
\item We isolate tokenization granularity as a single controlled variable in a
low-resource Korean--English recurrent NMT setting, correcting the confounds
present in the original comparison.
\item We characterize the interaction between subword vocabulary size and
training-data scale on translation quality, finding that BPE dominates
word-level tokenization across both scales while remaining largely insensitive
to vocabulary size, especially at full scale.
\item We separate \emph{representational} from \emph{optimization} explanations
of subword underperformance using stability diagnostics, providing evidence on
whether the observed failure is intrinsic.
\item We distill the findings into practical guidance for practitioners building
small recurrent NMT systems on limited parallel data.
\end{itemize}

% ======================================================================
\section{Background and Related Work}
% ======================================================================
\subsection{Attention-based recurrent NMT}
The encoder--decoder framework maps a source sequence to a fixed
representation and decodes a target sequence from it~\cite{sutskever2014,
cho2014}. Bahdanau et al.~\cite{bahdanau2015} removed the fixed-vector
bottleneck by letting the decoder attend to a weighted combination of all
encoder states, computed with an additive scoring function; Luong et
al.~\cite{luong2015} studied multiplicative alternatives. The attention
distribution doubles as a soft alignment, which makes recurrent attentional NMT
useful for interpretability studies even though Transformer architectures now
dominate at scale~\cite{vaswani2017}. We adopt the additive (Bahdanau)
formulation because it matches our pilot and because the resulting alignments
are convenient to visualize.

\subsection{Subword tokenization}
BPE~\cite{sennrich2016} iteratively merges the most frequent adjacent symbol
pairs, yielding a vocabulary of subword units that interpolates between
characters and words; SentencePiece~\cite{kudo2018spm} implements BPE (and a
unigram language-model variant) directly on raw text without language-specific
pre-tokenization, which is convenient for Korean. Subword regularization and the
unigram model~\cite{kudo2018subword} further improve robustness by sampling
segmentations. The number of merge operations, equivalently the vocabulary size,
is a free hyperparameter whose effect is often underexamined; Gowda and
May~\cite{gowda2020} argue that vocabulary size should be chosen with respect to
corpus statistics rather than copied from high-resource recipes, and report that
smaller vocabularies can be preferable when data are limited. Our granularity
sweep is in this spirit.

\subsection{Low-resource NMT and the limits of defaults}
Sennrich and Zhang~\cite{sennrich2019} show that low-resource NMT is acutely
sensitive to hyperparameters and that conclusions drawn at high resource do not
transfer unchanged; carefully tuned systems can substantially outperform
poorly tuned ones on the same data. This is directly relevant to our question:
an apparent representational disadvantage may instead reflect an untuned
optimizer interacting with longer subword sequences. By logging stability
diagnostics alongside quality, we aim to tell these explanations apart rather
than conflate them.

\subsection{Evaluation}
BLEU~\cite{papineni2002} remains the most reported MT metric, but its
comparability is undermined by inconsistent tokenization and preprocessing;
sacreBLEU~\cite{post2018} addresses this by computing a standardized score from
detokenized text. Because BLEU's $n$-gram matching is brittle for
morphologically rich targets, we also report chrF2~\cite{popovic2015}, a
character-$n$-gram F-score that credits partial morphological matches and
correlates better with human judgment on such languages.

% ======================================================================
\section{Methodology}
% ======================================================================
Our methodological contribution is not a new model but a corrected experimental
procedure. The pilot that motivates this work compared a word-level and a BPE
system that differed simultaneously in vocabulary size, sequence-length cap,
learning rate, and evaluation protocol, and judged them on a handful of
hand-selected sentences; such a comparison cannot attribute any observed
difference to tokenization. We improve on this in four concrete ways: (i)~we hold
every factor other than tokenization constant, including a vocabulary cap on the
word-level model; (ii)~we evaluate on a fixed, held-out test set rather than
cherry-picked examples; (iii)~we report standardized automatic metrics
(sacreBLEU, chrF2) instead of qualitative impressions; and (iv)~we log
optimization diagnostics so that quality differences can be separated from
training instability. The remainder of this section specifies the shared model,
the controlled design, and these diagnostics.

\subsection{Model}
All conditions share an identical architecture. The encoder is a single-layer
GRU over an embedding of dimension $d_e$; the decoder is a single-layer GRU
whose input at each step concatenates the previous-token embedding with a
context vector produced by additive attention over the encoder states. Given
encoder states $\mathbf{h}_1,\dots,\mathbf{h}_S$ and decoder state
$\mathbf{s}_{t-1}$, the attention energy and weights are
\begin{equation}
e_{t,i} = \mathbf{v}^\top \tanh(\mathbf{W}\mathbf{s}_{t-1} + \mathbf{U}\mathbf{h}_i),
\qquad
\alpha_{t,i} = \frac{\exp(e_{t,i})}{\sum_{j} \exp(e_{t,j})},
\end{equation}
and the context vector is
$\mathbf{c}_t = \sum_i \alpha_{t,i}\mathbf{h}_i$. Padding positions are masked
out of the softmax. The output projection maps the decoder state to target
vocabulary logits. Weights are initialized with Xavier uniform; embeddings use a
dedicated padding index.

\subsection{Controlled experimental design}
The study is a two-factor design.
\textbf{Factor~1 (tokenization granularity)} comprises a whitespace word-level
tokenizer with a capped vocabulary and BPE at vocabulary sizes
$\{2\text{k}, 4\text{k}, 8\text{k}\}$. We did not include a 1k BPE level:
SentencePiece requires the vocabulary to be at least as large as the set of
characters needed to cover the corpus, and for this Korean--English data that
floor (the Korean syllabary plus Latin characters and punctuation) exceeds 1k, so
a 1k model cannot represent the script and is rejected at training time. The
smallest admissible BPE setting is therefore 2k. We focused the granularity sweep
on the small-vocabulary regime ($\le$8k), which is the region of interest for
low-resource settings and the region where, as our results show, the
quality differences among BPE settings are concentrated.
\textbf{Factor~2 (data scale)} has two levels: 10k training pairs and the full
$\approx$71k-pair training set. The word-level vocabulary is capped at 16k to
keep it comparable to the BPE settings rather than allowing the
uncapped vocabulary used in the pilot. Every other quantity is held fixed across
all cells: embedding and hidden sizes, optimizer (AdamW), learning rate, weight
decay, gradient-clipping threshold, label smoothing, teacher-forcing ratio, batch
size, number of epochs, the sequence-length cap, the decoding procedure,
and---crucially---a single train/validation/test split that is drawn once and
reused everywhere. Tokenizers are fit on the training partition only. We use a
teacher-forcing ratio of 0.9; in preliminary runs a lower ratio caused the
decoder to collapse into degenerate repetition of high-frequency tokens under
greedy decoding, a known manifestation of exposure bias in small recurrent
decoders. Each cell is run with a fixed random seed; the released script supports
additional seeds for variance estimation, which we leave as a straightforward
extension. One full-scale cell (BPE-8k) did not complete due to an environment
interruption and is omitted; its absence does not affect any conclusion, as the
full-scale 2k and 4k cells already bracket its expected value.

\subsection{Disentangling representation from optimization}
To address the second half of the research question, each run logs three
optimization diagnostics in addition to its quality scores: the number of
batches discarded due to non-finite loss (a proxy for divergence), the peak
gradient norm observed during training, and the number of epochs required to
reach the best validation loss (a proxy for convergence speed). If subword
models trained to convergence match word-level quality, an apparent subword
disadvantage at fixed budget is an \emph{optimization} effect; if a gap persists
at convergence with stable training, it is better described as
\emph{representational}.

% ======================================================================
\section{Experimental Setup}
% ======================================================================
\subsection{Data}
We use the Korean--English parallel news corpus of
Park~\cite{park-corpus}.\footnote{Obtained from the public repository
\texttt{jungyeul/korean-parallel-corpora} (korean-english-news-v1).} The raw
release contains 94{,}123 sentence pairs. After deduplication and minimal,
tokenizer-agnostic normalization (whitespace collapsing only), 78{,}941 pairs
remain. We hold out 5\% of these as a test set and 5\% as a validation set with a
fixed random seed (42), leaving 71{,}047 training pairs; the 10k condition uses
the first 10{,}000 of these. Sentences are capped at 60 tokens, a generous bound
that is rarely active given the mean lengths reported in Table~\ref{tab:data}.

\begin{table}[t]
\caption{Corpus split statistics. Lengths are mean whitespace-token counts.}
\label{tab:data}
\centering
\begin{tabular}{lrrr}
\toprule
Split & Pairs & Mean KO len & Mean EN len \\
\midrule
Train      & 71{,}047 & 14.7 & 22.2 \\
Validation &  3{,}947 & 14.8 & 22.2 \\
Test       &  3{,}947 & 14.7 & 22.1 \\
\bottomrule
\end{tabular}
\end{table}

\subsection{Tokenizers}
Word-level tokenization splits on whitespace and caps the vocabulary at 16k,
mapping the remainder to an unknown token. BPE models are trained with
SentencePiece directly on the normalized text, with shared special tokens
(\texttt{pad}, \texttt{unk}, \texttt{bos}, \texttt{eos}) assigned identical
indices across all tokenizers so that the only difference between conditions is
the segmentation itself. SentencePiece is configured with a character coverage
of 0.9995.

\subsection{Training}
Models are optimized with AdamW (learning rate $5\times10^{-4}$, weight decay
$10^{-5}$) for 30 epochs, with gradient norm clipped to 1.0, label smoothing
0.1, and a fixed teacher-forcing ratio of 0.9. The learning rate is halved on
validation-loss plateau (patience 2). The loss is cross-entropy computed over all
non-padding target positions, and we select the checkpoint with the lowest
validation loss for testing. Embedding and hidden dimensions are 256 and 512,
respectively. Batch size is 128. Experiments run on a single
a single NVIDIA T4 GPU.

\subsection{Evaluation and decoding}
We decode the test set greedily and report corpus-level
sacreBLEU~\cite{post2018} and chrF2~\cite{popovic2015} against detokenized
references.

% ======================================================================
\section{Results}
% ======================================================================
\subsection{Quality across granularity and data scale}
Table~\ref{tab:main} reports test BLEU and chrF2 for every cell of the design.
The dominant effect is the tokenizer: BPE outperforms word-level tokenization by
a large margin at both data scales. At the full scale the best configuration was
BPE-4k with 4.19 BLEU / 23.88 chrF2, statistically indistinguishable from BPE-2k
(4.28 / 23.81), whereas word-level tokenization reached only 0.65 / 9.47. At the
10k scale the same ordering holds (BPE-2k 1.79 / 17.64 versus word-level
0.07 / 5.26). Word-level tokenization thus trails the best BPE setting by roughly
$14$ chrF2 points at full scale and by $12$ at 10k; the gap does not close as
data grows. This directly reverses the pilot observation, in which an uncapped
word-level model appeared to beat BPE: once vocabulary size, sequence-length cap,
learning rate, and the evaluation protocol are controlled, the apparent
word-level advantage disappears and BPE is uniformly preferable. We note that the
absolute BLEU scores are low throughout---an expected consequence of a
single-layer GRU trained on at most $\approx$71k pairs---but chrF2, which credits
partial morphological matches, separates the conditions cleanly and consistently.

\begin{table}[t]
\caption{Test BLEU / chrF2 (single seed, best-validation checkpoint). The
full-scale BPE-8k cell did not complete (environment interruption).}
\label{tab:main}
\centering
\begin{tabular}{l l cc cc}
\toprule
 & & \multicolumn{2}{c}{BLEU $\uparrow$} & \multicolumn{2}{c}{chrF2 $\uparrow$} \\
\cmidrule(lr){3-4}\cmidrule(lr){5-6}
Tokenizer & Vocab & 10k & Full & 10k & Full \\
\midrule
Word  & 16k & 0.07 & 0.65 & 5.26 & 9.47 \\
BPE   & 2k  & 1.79 & \textbf{4.28} & 17.64 & 23.81 \\
BPE   & 4k  & 1.60 & 4.19 & 15.83 & \textbf{23.88} \\
BPE   & 8k  & 1.30 & --   & 15.98 & -- \\
\bottomrule
\end{tabular}
\end{table}

\subsection{Effect of subword vocabulary size}
Figure~\ref{fig:vocab} plots chrF2 against BPE vocabulary size at each data
scale, with the word-level reference shown as a horizontal line. Within BPE, the
effect of vocabulary size is small and shrinks with data. At 10k the smallest
vocabulary (2k) is best (17.64) and performance is slightly lower and flat for
4k--8k (15.8--16.0); at full scale the 2k and 4k settings are essentially
identical (23.81 vs.\ 23.88). The curve is therefore nearly flat and far above
the word-level reference at both scales. This is consistent with the view that
once a subword vocabulary is large enough to represent the script compactly,
further enlargement yields diminishing returns, and that smaller vocabularies are
mildly preferable in the lowest-resource setting.

\begin{figure}[t]
\centering
\includegraphics[width=\columnwidth]{vocab_vs_chrf.png}
\caption{chrF2 versus BPE vocabulary size at each data scale; dotted lines are
the word-level references. BPE exceeds word-level by a wide margin everywhere,
while the BPE curves themselves are nearly flat.}
\label{fig:vocab}
\end{figure}

\subsection{Training stability}
Table~\ref{tab:std} reports the optimization diagnostics. With the loss computed
over all non-padding positions and a teacher-forcing ratio of 0.9, training was
stable in every cell: no run recorded a single non-finite-loss batch, and peak
gradient norms stayed near unity after clipping, with the sole exception of
BPE-8k at 10k (peak $\approx$10), the smallest-data/largest-vocabulary corner.
The clearest signal is in convergence speed rather than instability: word-level
training reached its best validation loss early (epoch 12 at full scale) and then
overfit, whereas the BPE models kept improving and reached their best checkpoints
later (epochs 20--30). In other words, the word-level model exhausts what it can
extract from the data sooner and at a worse optimum, while BPE continues to use
the additional data productively. The quality gap is therefore not attributable
to optimization failure of the BPE models---they train cleanly---but reflects a
genuine representational advantage of subword units for this language pair and
model class.

\begin{table}[t]
\caption{Optimization diagnostics (single seed). NaN batches, peak gradient norm,
and epoch at which the best validation loss was reached.}
\label{tab:std}
\centering
\begin{tabular}{l l rrr}
\toprule
Tokenizer & Vocab & Scale & NaN & Epochs-to-best \\
\midrule
Word & 16k & 10k  & 0 & 20 \\
Word & 16k & Full & 0 & 12 \\
BPE  & 2k  & 10k  & 0 & 25 \\
BPE  & 2k  & Full & 0 & 20 \\
BPE  & 8k  & 10k  & 0 & 30 \\
\bottomrule
\end{tabular}
\end{table}

\subsection{Qualitative output analysis}
Inspecting greedy-decoded outputs from the best full-scale BPE model corroborates
the quantitative picture and clarifies what the low BLEU scores do and do not
mean. The model produces fluent English fragments with correctly recovered
content words and named entities tied to the source---for example, rendering an
Iraq-and-diplomacy news sentence with the phrases ``the United States'' and
``President Roh Moo-hyun,'' and a military-incident sentence with ``the U.S.\
military,'' ``bombings,'' and ``government.'' These are not random high-frequency
tokens: they track the source topic, indicating that the attention mechanism is
routing source content to the decoder. Errors are of degradation rather than
collapse---occasional non-words, repeated clauses, and loss of long-range
coherence---consistent with the limited capacity of a single-layer GRU on a few
tens of thousands of pairs. By contrast, the word-level model and, in
preliminary low-teacher-forcing runs, all models, produced degenerate repetition
of the single most frequent token (``the the the\ldots''), which yields a near-zero
BLEU and a low chrF2. The gap between the conditions is thus not a metric artifact
but reflects a real difference between producing topical, partially correct
translations and producing none.

Korean and English differ in canonical word order (Korean is verb-final), so a
faithful model must learn non-monotonic alignments; the recovery of correctly
ordered subject--verb--object fragments in the better outputs is indirect
evidence that the attention layer learns some of this reordering, though we leave
a quantitative alignment-error analysis to future work.

% ======================================================================
\section{Discussion}
% ======================================================================
The controlled comparison lets us revisit the pilot's surprising result with the
confounds removed, and it overturns that result.

First, regarding \emph{when} word-level matches or beats BPE: in our controlled
setting, it never does. BPE is preferable at both data scales, and the
$\approx$14-point chrF2 gap at full scale is, if anything, larger than at 10k.
The pilot's apparent word-level advantage is best explained as an artifact of its
uncontrolled setup---an uncapped word vocabulary, a different sequence-length cap
and learning rate, and an evaluation on a handful of hand-picked sentences rather
than a held-out test set. This is consistent with prior cautions that conclusions
about tokenization are highly sensitive to experimental
configuration~\cite{gowda2020, sennrich2019}. Within BPE, the weak dependence on
vocabulary size echoes the finding that vocabulary should be matched to corpus
scale rather than imported from high-resource recipes~\cite{gowda2020}.

Second, regarding \emph{why}: the stability diagnostics show that the BPE models
are not winning by avoiding a pathology that afflicts word-level training. All
models trained stably (no non-finite batches, bounded gradients). The difference
is in how much each representation lets the model extract from the data:
word-level training saturated early and overfit (best epoch 12 at full scale),
whereas BPE continued to improve and converged later at a markedly lower loss.
The advantage is therefore representational rather than merely optimization-related:
subword units give a compact, less sparse vocabulary for an agglutinative
language, which a small recurrent decoder can exploit, while a word-level
vocabulary remains too sparse to learn from limited data. We separately observed
that teacher forcing was pivotal for obtaining any non-degenerate output at all:
at a ratio of 0.5 the decoder collapsed into repetition of frequent tokens under
greedy decoding, and only at 0.9 did the models produce content-bearing
translations---an exposure-bias effect that is itself worth flagging for
practitioners of small recurrent NMT.

We emphasize that these conclusions are scoped to small recurrent attentional
models on limited Korean--English data and should not be read as claims about
subword tokenization in general or about Transformer-scale systems.

% ======================================================================
\section{Limitations}
% ======================================================================
Several limitations bound our claims. The corpus is a single domain (news) and a
single language direction (Korean$\rightarrow$English); morphological effects may
differ in the reverse direction. We study a single-layer GRU architecture rather
than a Transformer, so capacity interacts with our findings. We use greedy
decoding; richer decoding (e.g.\ beam search) could shift absolute scores. BLEU
and chrF2 are imperfect proxies for adequacy and fluency, and we report no human
evaluation. We sweep BPE vocabulary size but hold the SentencePiece algorithm
fixed (BPE rather than the unigram model), and we do not use subword
regularization~\cite{kudo2018subword}. Finally, we report a single seed per cell;
repeating each configuration over multiple seeds would quantify variance and
tighten confidence around the smaller observed differences, and the released
script supports this directly.

% ======================================================================
\section{Conclusion}
% ======================================================================
We turned an anomalous pilot result---a word-level model appearing to outperform
a BPE model on low-resource Korean--English translation---into a controlled study
of tokenization granularity and training stability. By holding architecture,
optimization, and data splits fixed and varying only granularity and data scale,
we find that the pilot's result reverses: BPE outperforms word-level tokenization
at every data scale by a large and persistent margin (chrF2 23.8 vs.\ 9.5 at full
scale), while the choice of BPE vocabulary size has only a minor effect that
diminishes as data grows. Our stability analysis indicates that this advantage is
representational rather than an optimization artifact---all models trained
cleanly, but the word-level model saturated early and at a worse optimum. The
practical implications are concrete: practitioners building small recurrent NMT
systems on limited, morphologically rich data should prefer subword tokenization,
need not over-tune its vocabulary size within the small-vocabulary regime, and
should be wary of two pitfalls we encountered---exposure bias under low
teacher-forcing ratios, and uncontrolled comparisons that can invert the apparent
ranking of tokenizers. Future work includes extending the sweep to the unigram
SentencePiece model with subword regularization, repeating the design with a
Transformer backbone and multiple seeds, and adding a human adequacy evaluation.

% ======================================================================
\section*{Reproducibility}
% ======================================================================
All preprocessing, the fixed split, the model, the full ablation grid, and the
evaluation use the released script \texttt{nmt\_tokenization\_ablation.py},
available at \url{https://github.com/USERNAME/REPO} \PH{(replace with your repo
URL once published)}.

% ======================================================================
% ======================================================================
\begin{thebibliography}{00}

\bibitem{sutskever2014} I.~Sutskever, O.~Vinyals, and Q.~V.~Le,
``Sequence to sequence learning with neural networks,'' in
\emph{Advances in Neural Information Processing Systems (NeurIPS)}, 2014.

\bibitem{cho2014} K.~Cho \emph{et al.},
``Learning phrase representations using RNN encoder--decoder for statistical
machine translation,'' in \emph{Proc. EMNLP}, 2014.

\bibitem{bahdanau2015} D.~Bahdanau, K.~Cho, and Y.~Bengio,
``Neural machine translation by jointly learning to align and translate,'' in
\emph{Proc. ICLR}, 2015.

\bibitem{luong2015} M.-T.~Luong, H.~Pham, and C.~D.~Manning,
``Effective approaches to attention-based neural machine translation,'' in
\emph{Proc. EMNLP}, 2015.

\bibitem{vaswani2017} A.~Vaswani \emph{et al.},
``Attention is all you need,'' in
\emph{Advances in Neural Information Processing Systems (NeurIPS)}, 2017.

\bibitem{sennrich2016} R.~Sennrich, B.~Haddow, and A.~Birch,
``Neural machine translation of rare words with subword units,'' in
\emph{Proc. ACL}, 2016.

\bibitem{kudo2018spm} T.~Kudo and J.~Richardson,
``SentencePiece: A simple and language independent subword tokenizer and
detokenizer for neural text processing,'' in
\emph{Proc. EMNLP (System Demonstrations)}, 2018.

\bibitem{kudo2018subword} T.~Kudo,
``Subword regularization: Improving neural network translation models with
multiple subword candidates,'' in \emph{Proc. ACL}, 2018.

\bibitem{gowda2020} T.~Gowda and J.~May,
``Finding the optimal vocabulary size for neural machine translation,'' in
\emph{Findings of EMNLP}, 2020.

\bibitem{sennrich2019} R.~Sennrich and B.~Zhang,
``Revisiting low-resource neural machine translation: A case study,'' in
\emph{Proc. ACL}, 2019.

\bibitem{papineni2002} K.~Papineni, S.~Roukos, T.~Ward, and W.-J.~Zhu,
``BLEU: A method for automatic evaluation of machine translation,'' in
\emph{Proc. ACL}, 2002.

\bibitem{popovic2015} M.~Popovi\'c,
``chrF: Character $n$-gram F-score for automatic MT evaluation,'' in
\emph{Proc. Tenth Workshop on Statistical Machine Translation (WMT)}, 2015.

\bibitem{post2018} M.~Post,
``A call for clarity in reporting BLEU scores,'' in
\emph{Proc. Third Conf. on Machine Translation (WMT)}, 2018.

\bibitem{park-corpus} J.~Park, ``Korean parallel corpora: Korean--English news
(v1),'' GitHub repository, 2016. [Online]. Available:
\url{https://github.com/jungyeul/korean-parallel-corpora}

\end{thebibliography}

\end{document}
